\section{Synchoros Characterization Methodology}
In this section, we describe the proposed characterization methodology that leverages the synchoricity and its spatial invariance to achieve a scalable, one-time engineering effort to factor the impact of all wires in the design.
%allowing for an agile and post-route accurate estimate of average energy. 
%The latter is only possible when all wires in the CGRA fabrics are accounted for.
%to deliver post-route accurate average energy estimates. 

\textbf{Mapping and Configware Assumptions.} 
We assume that when an algorithm is mapped onto a customizable, heterogeneous CGRA fabric, the mapper (also referred to as a compiler/synthesis tool) produces the resulting configware.
This configware fully specifies the spatio-temporal distribution of operations executed by CGRA cells.
In practice, CGRA mappers commonly produce configware as a critical output for functional verification~\cite{ragheb2024cgra,wijerathne2022morpher, gobieski2021snafu, yang2022vesyla}. 

\subsection{Characterization Building Blocks}
The starting point for characterization is a set of abutment-ready SiLago blocks.
The blocks correspond to the CGRA PEs in a heterogeneous, customizable CGRA framework. 
Some CGRA frameworks, however, could be homogeneous and would constitute a subset of the assumed framework.
The blocks are synthesized as Design Rule Check (DRC) and timing-clean GDSII macros. 
%These blocks are also abutment-ready, with potential valid neighbors, to create a larger design that is also DRC and timing clean.  
The mapper determines the number and types of blocks required to implement a given mapping and also enumerates the operations for each block.

Each operation is bound to specific \textit{functional resources} within the corresponding SiLago block—for example,
ALUs to perform basic arithmetic and logic operations, or a register file (RF) for temporary data storage.
These resources get explicitly identified by the operations specified by the mapper.
Based on the CGRA design goal, the block may include more sophisticated functional resources, such as address generators (AGs), configured to access external memory in specific patterns. 
Overall, the number of functional resources is usually limited and can be profiled for characterization.

In a CGRA, Network-on-Chip (NoC) serves as the primary functional interconnect for data movement, while the clock tree (CT) serves as the infrastructural interconnect for clock distribution.
In synchoros CGRAs, each SiLago block includes a local NoC or \textit{LNoC interconnect resource} and a local CT or \textit{LCT interconnect resource}.
The final CGRA interconnect is composed by the abutment of local interconnect segments.
LNoC consists of switchboxes, buffers, and registers for pipelining. 
While LNoC is not explicitly identified by the configware, it is implied by the data-movement schedule dictated by the operations.
LCT is not identified or implied in operations but contributes to energy consumption and, hence, is also characterized.
% resources, such as switchbox traversal paths or inter-block routing ports, are implicit: they are not explicitly listed in the configware but are inferred from the data-movement schedule dictated by the mapper.
% There are two broad categories of operations that are characterized: 
% 1) Functional operations that can be potentially identified and implied by any configware. These categories of operations include DPU/ALU, Register File, AGUs, Sequencer, data transfer, etc. 
\subsection{Spatial-Context Aware Characterization}
A SiLago block is characterized in the context of its valid neighbours.
Valid neighbours include the span of PEs that are directly connected to each other. 
This enables to account for accurate drive strength and capacitive loads. 
For this purpose, a minimal SiLago design is generated by abutment. 
Consequently, the impact of all intra- and inter-block wires is inherently factored into the characterization. 
The set of valid neighbours is often much smaller than the total number of PEs in the design (see Figure~\ref{fig:spatial}).

%Table~\ref{table:wires} summarizes the types of interconnect at different levels of hierarchy in the CGRA fabric.
%and how their energy cost can be derived with agility and accuracy.
The energy cost of intra-component and inter-component wires is known because a SiLago block is characterized after the synthesis of the minimal design, after the locations of all transistors and wire segments within the block are fixed. 
As a result, the electrical impact of internal wires—both inside and between functional components—is included in the characterization of each resource.

The CGRA NoC and CT are composed by the abutment of pre-existing LNoC and LCT segments.
When two neighboring blocks abut, the corresponding wire segments align to create a larger, valid design without requiring re-synthesis or parasitics extraction.
Because each block is characterized in the context of its valid neighbors, inter-block electrical effects—such as variations in drive strength or capacitive loading—are inherently captured.

    \begin{figure}
        \centering
        \scriptsize
        \includesvg[scale=0.13]{figures/figures-spatial_context.drawio.svg}
        \caption{\small Example of two valid-neighbor configurations for synchoros characterization.}
        \label{fig:spatial}
    \end{figure}
%\input{tables/1_wires}

%Some operations are not directly identified by the configware but are implied, for instance, a switch box operation would be implied by a data transfer operation in the configware. 
% While characterizing these operations,  the impact of wires is factored in at three levels: a) all wires within the resources implementing these operations are factored in, b) intra-resource wires within the SiLago blocks are also factored in, and lastly, c) inter-SiLago block wires are also factored in when characterizing the operations. 
% This spatial-context-sensitive characterization is essential because, depending on the neighbors, the load and drive on the resources in the SiLago block would differ, resulting in different power and delay numbers.  
% The above strategy for characterizing operations while naturally factoring in wires at three levels is applied to all operations exported by resources in different types of SiLago blocks.
% 2) Infrastructural resources mainly refers to clock trees and local NoC wires and their buffers. These resources also contribute to significant power consumption and are also characterized and factored in during estimation, as elaborated in section xx. 

\begin{comment}
We identify two types of components within the SiLago block that are profiled for characterization based on the operations the block can perform:
\subsubsection{Functional components:} These components implement the core functionality of the PE, like 

\subsubsection{Interconnect components:} 



The local NoC segments within the block comprises the routers consisting of switchboxes, buffers, registers for pipelining. 
Local clock tree segment comprises of clock tree buffers.

\subsection{Spatial-Context-Aware Characterization}
SiLago block is characterized in the context of its valid neighbours.
Valid neighbours include the span of PEs that are directly connected to each other to communicate in one clock cycle.
This enables to account for accurate drive strength and capacitive loads. 
For this purpose, a minimal SiLago design is generated by abutment. 
Consequently, the impact of all intra- and inter-block wires are inherently factored into the characterisation.
It is one in typical neighbor-to-neighbor connections, with one neighbour in NESW direction.
\end{comment}
\subsection{Characterization Algorithm}
Within the minimal design, the individual block serves as the DUT(Design-Under-Test) for characterization. 
The block is simulated with the possible operations it supports. 
Configware information is used to profile the raw toggling activity of its resources during these operations. 
By analyzing the toggling patterns from simulations, we can determine power consumption for each operation and its associated hardware. 
The entire process results in a database of average power consumption for each block across all possible operating modes. 
Importantly, these power values account for all wires from the fabric level down to the inter-standard cell. 
The characterization algorithm for preparing the database of characterized operations is presented in Algorithm~\ref{alg:simconverge}.

To obtain accurate power values for every operation supported by the resources in a SiLago block, we employ an iterative characterization flow with a
convergence criteria across all contributing functional and interconnect components. 
For each block type $s \in S$, the algorithm enumerates all legal neighborhood configurations $N \subseteq \mathrm{ValidNeighbors}(s)$ and generates the corresponding SiLago design $D = \mathrm{Abutment}(s, N)$. 
For each operation $op \in \mathrm{Ops}(s)$ and each internal resource $r \in \mathrm{Resources}(op)$, the instantaneous power $P_{\mathrm{inst}}$ is repeatedly measured through post-layout parasitics back-annotated gate-level power simulation   of~$D$.
The running mean $P_{r,op}$ is updated as the online mean of all collected samples:
\begin{equation}
    P_{r,op}^{(k+1)}
    = \frac{k \, P_{r,op}^{(k)} + P_{\mathrm{inst}}^{(k+1)}}{k+1}.
\end{equation}


As gate-level simulations are time-consuming, we separate the characterization of static power consumption, which requires a single simulation iteration, from the dynamic power simulation, which requires multiple iterations. 
To minimize the number of iterations required to characterize dynamic power consumption, we have adopted a convergence-based strategy, as shown in Figure~\ref {fig:convergence}. 
To ensure that the characterization does not introduce bias, random stimuli are applied, and a running average is maintained. 
\begin{comment}
Gate-level simulations are time-consuming. 
The gate count of a single block can be at least tens to hundreds of thousands, requiring hours of simulation effort.
It is in our interest to minimize the number of simulations while still performing sufficient simulations for good confidence in characterization.
Static power, independent of switching activity, can be captured accurately in a single simulation. 
Dynamic power, on the other hand, depends heavily on the switching behavior of the circuit and requires multiple simulations. 
CMOS circuits consume the highest power when switching.
To ensure that the characterized power values are representative of the various switching that may arise in the circuit, test vectors are applied with a series of randomized inputs that are repeated until the values stabilize. 
To reduce the simulation overhead, we adopt a convergence-based approach.
\end{comment}
\input{algorithms/algorithm1}
After a minimum number of simulations is performed, the error between consecutive running averages of power is tracked over a window of simulations. 
If the errors in all of the simulations within the window are below the convergence threshold, the simulations are terminated. 
The threshold and the simulation window are user-defined.

To ensure stability across varying inputs, convergence is not declared immediately.
 Instead, after an initial settling period of $n$ iterations, the algorithm tracks the absolute deviation between the latest sample and the running mean,
\begin{equation}
    e^{(k)} = \left|P_{\mathrm{inst}}^{(k)} - P_{r,op}^{(k)}\right|,
\end{equation}
and stores it in a sliding window of size $w$.
Convergence is reached only when all errors in the window satisfy
\begin{equation}
    e^{(j)} < \epsilon
    \qquad
    \forall j \in \{k-w+1, \ldots, k\}.
\end{equation}
This condition guarantees that the resource-level power estimate has stabilized and that further simulations would not significantly affect the final average. The resulting converged value $P_{r,op}$ is recorded as the characterized power of resource $r$ when executing operation $op$ in the given neighborhood configuration.
\begin{figure}[t]
    \centering
    \scriptsize
\input{figures/3_convergence}
\caption{\small Iterative convergence of the running-mean power estimate for a single resource under one operation.}
\label{fig:convergence}
\end{figure}

All converged power values are finally stored in a \textit{characterization database} that consolidates the results across all blocks and its resources, operations, and valid neighborhood configurations. 
Each entry in the database is indexed by a tuple,
\[
(s,\; op,\; r,\; N)
\]
corresponding respectively to the SiLago block type, the operation being executed, the contributing resource, and the neighborhood context in which characterization was performed. 
Because each block type supports only a modest number of operations with its limited functional and interconnect resources in the context of its valid neighbours, the dimensionality of this database remains small — typically on the order of tens to a few hundreds of entries for a heterogeneous CGRA fabric. 
As a result, the database is compact and scalable: adding a new block type increases its size linearly without requiring global re-characterization. 
For estimation, the required values can be retrieved by directly querying the database. 
%In this way, each characterized value serves as one reusable ``atom'' in a larger, once-engineered catalogue that captures the power behavior of the entire fabric under all relevant spatial contexts.
